\subsection{The first-passage gate: Tao's Proposition~1.11}

The controlling statement is Proposition~1.11 at v7 source
lines~316--335, and the deduction from Propositions~1.9 and~1.14 occupies
v7 lines~649--925 \cite{Tao2026v7}.  The corresponding comparison loci are
v5 lines~643--919 and journal physical pages~18--27.  The three
manifestations have the same substantive modular, logarithmic-mass, and
integer-window calculations.  ArXiv v7 corrects the v5/journal interchange
of passage time and passage location in the converse implication (v5
lines~789--791; journal physical page~23; v7 lines~795--797).  No conclusion
below is transferred silently between versions.

\begin{sourceissue}[local defects at the Proposition~1.11 edge]
\label{sourceissue:Tao-Prop111-edge}
The opening heuristic at v7 lines~651--652 writes
\(\operatorname{Pass}_x\) where the quantity being approximated is the time
\(T_x\), and its numerator requires \(\log(N_y/x)\).  Line~689 has
\(1.9n\) where the fixed descent time is \(1.9n_0\).  The periodic-orbit
aside at line~699 requires a representative of \(-\mathbf m\) modulo the
least positive period, or a period multiple at least \(\mathbf m\); that
aside is not used here.

The additive endpoint deletion printed at lines~759--762 does not produce a
time margin of order \((\log x)^{0.8}\).  It has already been replaced by the
multiplicative deletion in
Proposition~\ref{prop:Tao-v7-endpoint-buffer}.  The unsigned lower comparison,
the event equality, the undefined \(m\) in the tuple length, the fixed-
\(M\) affine equality, the two previous-iterate coefficients, the malformed
progression estimate, the nonuniform weakened \(c_n\)-majorant, and the
iid-versus-actual Chernoff citation have already been repaired in
Proposition~\ref{prop:Tao-v7-Prop52-cn-repair}.  This subsection supplies the
remaining modular law, descent exponent, affine inverse, integer count,
normalization cancellation, and arbitrary-event-to-total-variation arrows.
\end{sourceissue}

Throughout this subsection retain the source constants and floors
\begin{equation}
\label{eq:Tao-Prop111-parameters}
 \alpha=\frac{1001}{1000},\qquad
 L=\log x,\qquad \lambda=\log(4/3),\qquad
 n_0=\left\lfloor\frac{L}{10\log2}\right\rfloor,
 \qquad m_0=\left\lfloor\frac{\alpha-1}{100}L\right\rfloor.
\end{equation}
For \(j\in\{1,2\}\), put \(y_j=x^{\alpha^j}\) and
\begin{equation}
\label{eq:Tao-Prop111-support-normalizer}
 \mathcal R_{y_j}=(2\mathbb N+1)\cap[y_j,y_j^\alpha],
 \qquad
 H_{y_j}=\sum_{N\in\mathcal R_{y_j}}\frac1N,
 \qquad
 \mathbb P(\mathbf N_{y_j}=N)=\frac{1}{H_{y_j}N}.
\end{equation}
Real endpoints cause no ambiguity: the intersection is with the integer set
\(2\mathbb N+1\).

\begin{lemma}[progression mass with real endpoints]
\label{lem:Tao-Prop111-progression-mass}
Let \(1\le A<B\), let \(q\ge2\) be even, put
\(J=\log(B/A)\), and fix an odd residue \(r\pmod q\).  Then
\begin{equation}
\label{eq:Tao-Prop111-one-progression}
 \left|
  \sum_{\substack{A\le N\le B\\N\equiv r\ (\mathrm{mod}\ q)}}\frac1N
  -\frac Jq
 \right|\le\frac1A.
\end{equation}
If \(W_r\) denotes the sum in
\eqref{eq:Tao-Prop111-one-progression} and
\(H=\sum_{r\ {\rm odd}\ (q)}W_r\), then
\begin{equation}
\label{eq:Tao-Prop111-total-progression-error}
 \sum_{r\ {\rm odd}\ (q)}\left|W_r-\frac Jq\right|
 \le\frac q{2A},
 \qquad
 \left|H-\frac J2\right|\le\frac q{2A}.
\end{equation}
\end{lemma}

\begin{proof}
If the progression contributes points
\(t_0<t_1<\cdots<t_s\), consecutive points differ by \(q\).  Monotonicity
of \(1/t\) gives
\[
 \sum_{i=0}^s\frac1{t_i}
 \le \frac1{t_0}+\frac1q\int_{t_0}^{t_s}\frac{dt}{t}
 \le \frac1A+\frac Jq.
\]
It also gives
\[
 \sum_{i=0}^s\frac1{t_i}
 \ge \frac1q\int_{t_0}^{t_s+q}\frac{dt}{t}
 \ge \frac Jq-\frac1q\log\frac{t_0}{A}
 \ge \frac Jq-\frac1A,
\]
because \(0\le t_0-A<q\) and
\(\log(t_0/A)\le(t_0-A)/A\).  If the progression contributes no point,
then \(B-A<q\), whence \(J/q<(B-A)/(qA)<1/A\).  This proves
\eqref{eq:Tao-Prop111-one-progression}.  There are exactly \(q/2\) odd
residue classes.  Summing and using the triangle inequality proves
\eqref{eq:Tao-Prop111-total-progression-error}.
\end{proof}

\begin{proposition}[exact modular input to Proposition~1.9]
\label{prop:Tao-Prop111-modular-input}
Let \(q=2^{3n_0}\).  Uniformly for \(j=1,2\),
\begin{equation}
\label{eq:Tao-Prop111-modular-law}
 d_{\mathrm{TV}}\!\left(
  \mathbf N_{y_j}\bmod q,
  \operatorname{Unif}((2\mathbb Z+1)/q\mathbb Z)
 \right)\ll q^{-1}=2^{-3n_0},
\end{equation}
where total variation is the source's unhalved \(\ell^1\) distance.  Hence
Proposition~\ref{prop:Tao-v7-Prop19-valuation-law}, with
\[
 n=n_0,\qquad n'=3n_0,\qquad c_0=1,
\]
gives
\begin{equation}
\label{eq:Tao-Prop111-valuations}
 d_{\mathrm{TV}}\!\left(
  \vec a^{(n_0)}(\mathbf N_{y_j}),
  \operatorname{Geom}(2)^{n_0}
 \right)\ll x^{-c_2}
\end{equation}
for one absolute \(c_2>0\).
\end{proposition}

\begin{proof}
Set \(A=y_j\), \(B=y_j^\alpha\), and
\(J=(\alpha-1)\log y_j\).  For every odd \(r\pmod q\), let \(W_r\) be
the sum in Lemma~\ref{lem:Tao-Prop111-progression-mass}.  For sufficiently
large \(x\), \(AJ\ge2q\), and therefore
\(H_{y_j}\ge J/4\).  Since uniform measure assigns mass \(2/q\), not
\(1/q\), to each of the \(q/2\) odd classes,
\begin{align}
 d_{\mathrm{TV}}
 &=\sum_{r\ {\rm odd}\ (q)}
   \left|\frac{W_r}{H_{y_j}}-\frac2q\right| \notag\\
 &\le \frac1{H_{y_j}}
 \left(
  \sum_{r\ {\rm odd}\ (q)}\left|W_r-\frac Jq\right|
  +\left|H_{y_j}-\frac J2\right|
 \right)
 \le\frac{2q}{AJ-q}\le\frac{4q}{AJ}.
\label{eq:Tao-Prop111-modular-L1-bound}
\end{align}
The floor in \(n_0\) is retained exactly:
\[
 \frac{x^{3/10}}8\le q\le x^{3/10}.
\]
Consequently
\[
 \frac{4q}{AJ}
 \le q^{-1}
 \frac{4x^{0.6-\alpha}}
      {(\alpha-1)\alpha L}
 =o(q^{-1}),
\]
uniformly for \(j=1,2\).  This proves
\eqref{eq:Tao-Prop111-modular-law} with room to spare.  The size hypothesis
in Proposition~\ref{prop:Tao-v7-Prop19-valuation-law} is the exact equality
\(3n_0=(2+c_0)n_0\).  Finally,
\[
 2^{-c_1n_0}\le2^{c_1}x^{-c_1/10},
\]
so its conclusion has the form \eqref{eq:Tao-Prop111-valuations} despite the
floor.
\end{proof}

\begin{proposition}[the finite-time descent estimate]
\label{prop:Tao-Prop111-finite-descent}
There is an absolute \(c_3>0\) such that, uniformly for \(j=1,2\),
\begin{equation}
\label{eq:Tao-Prop111-failure}
 \mathbb P\bigl(T_x(\mathbf N_{y_j})=+\infty\bigr)\ll x^{-c_3}.
\end{equation}
More precisely, outside an event of probability \(O(x^{-c_3})\),
\(T_x(\mathbf N_{y_j})\le n_0\).
\end{proposition}

\begin{proof}
For iid \(G_1,\ldots,G_{n_0}\sim\operatorname{Geom}(2)\), the lower-tail
Chernoff bound applied to
\[
 \mathbb E e^{-tG_1}=\frac{e^{-t}}{2-e^{-t}}
\]
gives
\(\mathbb P(G_1+\cdots+G_{n_0}\le1.9n_0)\le e^{-\gamma n_0}\)
for an absolute \(\gamma>0\).  For example, take
\(e^{-t}=18/19\); the base is
\(0.9(19/18)^{1.9}<1\).  Combining this with
\eqref{eq:Tao-Prop111-valuations} gives
\begin{equation}
\label{eq:Tao-Prop111-low-sum}
 \mathbb P\!\left(
  |\vec a^{(n_0)}(\mathbf N_{y_j})|\le1.9n_0
 \right)\ll x^{-c_3}
\end{equation}
after decreasing \(c_3\) if necessary.

On the complementary event, the exact affine formula and
\(0\le F_{n_0}\le3^{n_0}\) give
\begin{equation}
\label{eq:Tao-Prop111-descent-bound}
 \operatorname{Syr}^{n_0}(\mathbf N_{y_j})
 \le 3^{n_0}2^{-1.9n_0}x^{\alpha^3}+3^{n_0}.
\end{equation}
The numerical exponent check is strict, not heuristic.  Put
\[
 \beta_1=\alpha^3+
  \frac{\log3-1.9\log2}{10\log2},
 \qquad
 \beta_2=\frac{\log3}{10\log2}.
\]
Rational interval arithmetic gives
\begin{equation}
\label{eq:Tao-Prop111-exponents}
 \beta_1<0.972<1,
 \qquad \beta_2<0.159<1.
\end{equation}
One elementary certificate for these decimal inequalities uses
\[
 \log u=2\sum_{k=0}^{M}\frac{z^{2k+1}}{2k+1}+R_M,
 \quad z=\frac{u-1}{u+1},
 \quad 0<R_M<\frac{2z^{2M+3}}{(2M+3)(1-z^2)},
\]
with \(u=2,3\), \(z=1/3,1/2\), and \(M=20\); every endpoint is rational.
Since
\(n_0\ge L/(10\log2)-1\) and
\(\log3-1.9\log2<0\), the first term on the right side of
\eqref{eq:Tao-Prop111-descent-bound} is at most
\[
 e^{1.9\log2-\log3}x^{\beta_1}.
\]
The positive coefficient in the second term instead uses
\(n_0\le L/(10\log2)\), and that term is at most \(x^{\beta_2}\).
Thus
\(e^{1.9\log2-\log3}x^{\beta_1}+x^{\beta_2}<x\) for all sufficiently
large \(x\).
Thus passage occurs by time \(n_0\), and
\eqref{eq:Tao-Prop111-low-sum} proves the claim.
\end{proof}

For the passage-location part, let \(E\subseteq(2\mathbb N+1)\cap[1,x]\).
For either choice of \(y=y_j\), retain the source interval
\begin{equation}
\label{eq:Tao-Prop111-Iy}
 I_y=\left[
  \frac{\log(y/x)}\lambda+L^{4/5},
  \frac{\log(y^\alpha/x)}\lambda-L^{4/5}
 \right]
\end{equation}
and the source set, independent of \(y\),
\begin{equation}
\label{eq:Tao-Prop111-Eprime}
\begin{split}
 E'=\{M\in2\mathbb N+1:\;&T_x(M)=m_0,
 \ \operatorname{Pass}_x(M)\in E,\\
 &e^{-L^{7/10}}(4/3)^{m_0}x
 \le M\le
 e^{L^{7/10}}(4/3)^{m_0}x\}.
\end{split}
\end{equation}
For an integer \(n\in I_y\), put \(r=n-m_0\).  For sufficiently large
\(x\), every such \(n\) satisfies
\begin{equation}
\label{eq:Tao-Prop111-range}
 m_0\le r\le n_0-m_0,
 \qquad 1\le m_0\le r.
\end{equation}

\begin{proposition}[the singleton affine inverse]
\label{prop:Tao-Prop111-affine-inverse}
Let \(n\in I_y\cap\mathbb Z\), \(r=n-m_0\), and
\(\vec a=(a_1,\ldots,a_r)\in\mathcal A^{(r)}\).  Put
\[
 A_0=0,
 \qquad A_k=a_1+\cdots+a_k,
 \qquad A=A_r,
 \qquad
 C_r(\vec a)=\sum_{i=1}^r3^{r-i}2^{A_{i-1}},
 \qquad F_r(\vec a)=2^{-A}C_r(\vec a).
\]
For \(M\in E'\), the equation
\(\operatorname{Aff}_{\vec a}(N)=M\) has at most one solution, namely
\begin{equation}
\label{eq:Tao-Prop111-Phi}
 \Phi_{r}(\vec a,M)
 =\frac{2^AM-C_r(\vec a)}{3^r}
 =2^A\frac{M-F_r(\vec a)}{3^r}.
\end{equation}
It has a solution in \(\mathcal R_y\) if and only if
\begin{equation}
\label{eq:Tao-Prop111-congruence}
 M\equiv F_r(\vec a)\pmod{3^r}
\end{equation}
in \(\mathbb Z[1/2]/3^r\mathbb Z[1/2]\).  When this congruence holds,
\(\Phi_r(\vec a,M)\) is a positive odd integer in \([y,y^\alpha]\),
\[
 \vec a^{(r)}(\Phi_r(\vec a,M))=\vec a,
 \qquad
 \operatorname{Syr}^r(\Phi_r(\vec a,M))=M,
\]
and its logarithmic point mass is
\begin{equation}
\label{eq:Tao-Prop111-point-mass}
 \mathbb P\!\left(
  \operatorname{Aff}_{\vec a}(\mathbf N_y)=M
 \right)
 =\frac{1+O(x^{-c_4})}{\frac{\alpha-1}{2}\log y}
  \frac{2^{-A}3^r}{M}
\end{equation}
for one absolute \(c_4>0\), uniformly in all displayed data.
\end{proposition}

\begin{proof}
The affine identity
\[
 \operatorname{Aff}_{\vec a}(N)=\frac{3^rN+C_r(\vec a)}{2^A}
\]
proves uniqueness and forces \eqref{eq:Tao-Prop111-Phi}.  Congruence
\eqref{eq:Tao-Prop111-congruence} is necessary.  Conversely, it makes
\eqref{eq:Tao-Prop111-Phi} an integer: \(M-F_r(\vec a)\) belongs to
\(3^r\mathbb Z[1/2]\), and its denominator divides \(2^A\).  The integer is
odd because \(C_r(\vec a)\) is odd, \(2^AM\) is even, and \(3^r\) is odd.

It remains to prove positivity and the interval bounds, rather than infer
them from presentation.  Since \(r\le n_0\),
\[
 0\le F_r(\vec a)\le3^r\le x^{0.159}
\]
by \eqref{eq:Tao-Prop111-exponents}.  From \eqref{eq:Tao-Prop111-Eprime},
\(M\ge e^{-L^{7/10}}x\).  Since
\(L^{7/10}\le(141/1000)L\) for all sufficiently large \(L\), the preceding
two bounds give
\begin{equation}
\label{eq:Tao-Prop111-offset-ratio}
 0\le\frac{F_r(\vec a)}M\le x^{-7/10}<\frac12.
\end{equation}
Write \(A=2r+\delta\) and
\(M=x(4/3)^{m_0}e^\eta\).  Membership in
\(\mathcal A^{(r)}\) and \(E'\) gives
\(|\delta|<L^{3/5}\) and \(|\eta|\le L^{7/10}\).  Hence
\begin{equation}
\label{eq:Tao-Prop111-Phi-log}
 \log\Phi_r(\vec a,M)
 =L+n\lambda+\delta\log2+\eta+
   \log\left(1-\frac{F_r(\vec a)}M\right).
\end{equation}
The final three terms have absolute value at most \(2L^{7/10}\) for all
sufficiently large \(x\).  The endpoints of \(I_y\) give margins
\(\lambda L^{4/5}\), and
\(2L^{7/10}<\lambda L^{4/5}\) eventually.  Thus
\(\log y\le\log\Phi_r\le\log y^\alpha\), proving both positivity and the
claimed range.  Substitution proves
\(\operatorname{Aff}_{\vec a}(\Phi_r)=M\); the exact valuation-prefix lemma,
Lemma~\ref{lem:Tao-Prop19-residue-fibre}, then proves the valuation and
iterate identities.

Finally, Lemma~\ref{lem:Tao-Prop111-progression-mass} with modulus \(2\)
gives
\begin{equation}
\label{eq:Tao-Prop111-harmonic-normalizer}
 \left|H_y-\frac{\alpha-1}{2}\log y\right|\le\frac1y,
 \qquad
 H_y=\frac{\alpha-1}{2}\log y\,(1+O(x^{-1})).
\end{equation}
Using \eqref{eq:Tao-Prop111-Phi} exactly,
\[
 \frac1{\Phi_r(\vec a,M)}
 =\frac{2^{-A}3^r}{M-F_r(\vec a)}.
\]
Equations~\eqref{eq:Tao-Prop111-offset-ratio} and
\eqref{eq:Tao-Prop111-harmonic-normalizer} now give
\eqref{eq:Tao-Prop111-point-mass}.
\end{proof}

The last proposition is an exact morphism, not merely a pointwise formula.
Define
\[
 \mathcal D_{r,y}=
 \{(\vec a,M)\in\mathcal A^{(r)}\times E':
       M\equiv F_r(\vec a)\pmod{3^r}\}
\]
and
\[
 \mathcal X_{r,y}=
 \{N\in\mathcal R_y:
   \vec a^{(r)}(N)\in\mathcal A^{(r)},\
   \operatorname{Syr}^r(N)\in E'\}.
\]
Then
\begin{equation}
\label{eq:Tao-Prop111-bijection}
 \Theta_{r,y}:\mathcal X_{r,y}\longrightarrow\mathcal D_{r,y},
 \qquad
 N\longmapsto
 \bigl(\vec a^{(r)}(N),\operatorname{Syr}^r(N)\bigr)
\end{equation}
is a bijection with inverse \(\Phi_r\).  Both compositions are identities;
every fibre is a singleton; its set-theoretic kernel pair is the diagonal;
and no coordinate or higher binary digit is discarded.  This is the exact
map underlying source lines~808--827.

\begin{lemma}[the exact integer time-window count]
\label{lem:Tao-Prop111-window-count}
For \(y_j=x^{\alpha^j}\), let the endpoints of \(I_{y_j}\) be
\[
 a_j=\frac{\alpha^j-1}{\lambda}L+L^{4/5},
 \qquad
 b_j=\frac{\alpha^{j+1}-1}{\lambda}L-L^{4/5}.
\]
For sufficiently large \(x\), the interval is nonempty and
\begin{equation}
\label{eq:Tao-Prop111-exact-count}
\begin{split}
 \#(I_{y_j}\cap\mathbb Z)
 &=\lfloor b_j\rfloor-\lceil a_j\rceil+1\\
 &=\frac{\alpha-1}{\lambda}\log y_j-2L^{4/5}+\varepsilon_j,
 \qquad -1<\varepsilon_j\le1.
\end{split}
\end{equation}
Consequently
\begin{equation}
\label{eq:Tao-Prop111-normalization-cancellation}
 \frac{\#(I_{y_j}\cap\mathbb Z)}
      {\frac{\alpha-1}{2}\log y_j}
 =\frac2\lambda
  -\frac4{(\alpha-1)\alpha^j}L^{-1/5}
  +O(L^{-1})
 =\frac2\lambda+O(L^{-1/5}).
\end{equation}
The leading term is independent of \(j\); the displayed finite-
\(x\) buffer correction is not.
\end{lemma}

\begin{proof}
The first line is the exact count of integer points in a closed real
interval.  Subtracting the endpoints gives the second line with
\[
 \varepsilon_j=(\lfloor b_j\rfloor-b_j)
 +(a_j-\lceil a_j\rceil)+1,
 \qquad -1<\varepsilon_j\le1.
\]
Since \(\log y_j=\alpha^jL\), division proves
\eqref{eq:Tao-Prop111-normalization-cancellation}.  It also shows why the
source's variable exponent \(c\) must be decreased to at most \(1/5\) at
this step.  The lower endpoint satisfies \(a_j\ge2m_0\), and
\[
 \frac{\alpha^3-1}{\lambda}<0.010439
 <0.144269<\frac1{10\log2};
\]
therefore \(I_y\subseteq[2m_0,n_0]\) for all sufficiently large \(x\),
which also proves \eqref{eq:Tao-Prop111-range} with every floor retained.
\end{proof}

\begin{theorem}[Tao Proposition~1.11, exact reconstructed edge]
\label{thm:Tao-v7-Prop111-first-passage}
Let \(\alpha=1.001\), and let \(\mathbf N_y\) have the logarithmic law on
the odd integers in \([y,y^\alpha]\).  There is an absolute \(c>0\) such
that, for all sufficiently large \(x\),
\begin{equation}
\label{eq:Tao-Prop111-theorem-failure}
 \mathbb P\bigl(T_x(\mathbf N_y)=+\infty\bigr)\ll x^{-c}
 \qquad (y=x^\alpha,x^{\alpha^2})
\end{equation}
and
\begin{equation}
\label{eq:Tao-Prop111-theorem-TV}
 d_{\mathrm{TV}}\!\left(
  \operatorname{Pass}_x(\mathbf N_{x^\alpha}),
  \operatorname{Pass}_x(\mathbf N_{x^{\alpha^2}})
 \right)\ll L^{-c}
\end{equation}
in the source's unhalved convention.
\end{theorem}

\begin{proof}
Equation~\eqref{eq:Tao-Prop111-theorem-failure} is
Proposition~\ref{prop:Tao-Prop111-finite-descent}.  Fix an arbitrary
\(E\subseteq(2\mathbb N+1)\cap[1,x]\), form \(E'\) by
\eqref{eq:Tao-Prop111-Eprime}, and define
\begin{equation}
\label{eq:Tao-Prop111-ZE}
 Z_E=\sum_{M\in E'}\frac{3^{m_0}}M
 \mathbb P\!\left(
  \operatorname{Syrac}(\mathbb Z/3^{m_0}\mathbb Z)
  \equiv M\pmod{3^{m_0}}
 \right).
\end{equation}
This quantity depends on \(x\) and \(E\), but not on whether
\(y=x^\alpha\) or \(y=x^{\alpha^2}\).

Proposition~\ref{prop:Tao-v7-endpoint-buffer} places the first-passage time
in \(I_y\) outside an \(O(L^{-1/5})\) event.  The exact event equality and
signed pre-passage estimate in
Proposition~\ref{prop:Tao-v7-Prop52-cn-repair}, together with the bijection
\eqref{eq:Tao-Prop111-bijection} and point mass
\eqref{eq:Tao-Prop111-point-mass}, give
\begin{equation}
\label{eq:Tao-Prop111-pre-mixing}
\begin{split}
 \mathbb P(\operatorname{Pass}_x(\mathbf N_y)\in E)
 ={}&\frac{1+O(x^{-c_4})}
          {\frac{\alpha-1}{2}\log y}
 \sum_{n\in I_y\cap\mathbb Z} A_{n,E}
 +O(L^{-c_5}),
\end{split}
\end{equation}
where, with \(r=n-m_0\),
\[
 A_{n,E}=3^r
 \sum_{\vec a\in\mathcal A^{(r)}}2^{-|\vec a|}
 \sum_{\substack{M\in E'\\M\equiv F_r(\vec a)\ (3^r)}}\frac1M.
\]
All sums are over their displayed discrete sets; no conditional law or
coupling is being inserted.

The uniform \(c_n\)-bound in
Proposition~\ref{prop:Tao-v7-Prop52-cn-repair} permits deletion of the bad
iid path tube.  Proposition~\ref{prop:Tao-v7-Fourier-to-mixing}, whose
dependency branch is closed in
Corollary~\ref{cor:Tao-v7-fine-scale-branch}, then replaces the law at level
\(r\) by the exact uniform lift of its level-\(m_0\) projection.  The finite
pairing is
\begin{equation}
\label{eq:Tao-Prop111-exact-collapse}
\begin{split}
 &\sum_{X\bmod3^r}
 \left(3^r\!\sum_{\substack{M\in E'\\M\equiv X\ (3^r)}}\frac1M\right)
 3^{m_0-r}
 \mathbb P\!\left(
  \operatorname{Syrac}(\mathbb Z/3^{m_0}\mathbb Z)
  \equiv X\pmod{3^{m_0}}
 \right)\\
 &\hspace{35mm}=Z_E.
\end{split}
\end{equation}
Indeed, for each \(M\) there is exactly one residue
\(X=M\pmod{3^r}\), and \(3^r3^{m_0-r}=3^{m_0}\).  Thus, choosing the
arbitrary mixing exponent large enough before summing over the
\(O(L)\) possible times,
\begin{equation}
\label{eq:Tao-Prop111-uniform-An}
 A_{n,E}=Z_E+O(L^{-c_6})
\end{equation}
uniformly for both \(y\)'s and all \(n\in I_y\cap\mathbb Z\).  Moreover
\(A_{n,E}=O(1)\) by the uniform \(c_n\)-bound, so
\eqref{eq:Tao-Prop111-uniform-An} gives \(Z_E=O(1)\) without circular use
of the probability conclusion.

Substitute \eqref{eq:Tao-Prop111-uniform-An} into
\eqref{eq:Tao-Prop111-pre-mixing}.  The \(O(L)\) summation factor is
cancelled by the denominator of size \(L\); after decreasing the common
positive exponent,
\[
 \mathbb P(\operatorname{Pass}_x(\mathbf N_y)\in E)
 =\frac{\#(I_y\cap\mathbb Z)}
        {\frac{\alpha-1}{2}\log y}Z_E+O(L^{-c}).
\]
Lemma~\ref{lem:Tao-Prop111-window-count} therefore yields the uniform,
\(y\)-independent main term
\begin{equation}
\label{eq:Tao-Prop111-arbitrary-event}
 \mathbb P(\operatorname{Pass}_x(\mathbf N_y)\in E)
 =\frac{2}{\lambda}Z_E+O(L^{-c}).
\end{equation}
Taking the difference of \eqref{eq:Tao-Prop111-arbitrary-event} for the two
values of \(y\), and then the supremum over all subsets \(E\) of the common
state space \((2\mathbb N+1)\cap[1,x]\), gives an \(O(L^{-c})\) event
distance.  The convention \(\operatorname{Pass}_x(N)=1\) when
\(T_x(N)=+\infty\) keeps both laws in this same state space.  Finally the
source inequality
\[
 d_{\mathrm{TV}}(P,Q)
 \le2\sup_E|P(E)-Q(E)|
\]
proves \eqref{eq:Tao-Prop111-theorem-TV} in the unhalved convention.
\end{proof}

\begin{nonclaim}
Theorem~\ref{thm:Tao-v7-Prop111-first-passage} proves exactly the
first-passage stabilisation proposition and closes the edge
\[
 \text{Proposition~1.9}+\text{Proposition~1.14}
 \Longrightarrow\text{Proposition~1.11}.
\]
It does not assert exact residue uniformity, independence of actual valuation
coordinates, arbitrary-\(y\) uniformity, natural-density transport, a coupling
of the two passage locations, convergence to \(1\), exclusion of divergent or
periodic orbits, or a fixed universal orbit bound.  The finite-
\(x\) time-window correction depends on \(y\), even though its leading
normalization does not.  The exact later implications to Theorem~3.1,
Theorem~1.6, and Theorem~1.3, including their source defects and repaired
morphisms, are proved in the next subsection; they had not yet been proved at
the checkpoint represented by this subsection.
\end{nonclaim}
